\section*{Capítulo \ref{ch:b1:radiometric-methods} --- Métodos de datação radiométrica}
% values: calc/ch25.py

\begin{solution}{exo:b1:radiometric-methods:1}
(a)~$D^*/P = 2^n - 1 = 3$ dá $n = 2$: duas meias-vidas. (b)~Rubídio-87:
$e^{\lambda t} - 1 = e^{\num{1.3972e-11} \times \num{3.10e8}} - 1 = 0.004341$;
urânio-235: $e^{\num{9.8485e-10} \times \num{3.10e8}} - 1 = 0.3570$. A razão do urânio é
cerca de 80 vezes maior: um dado erro absoluto de medida é uma parte bem menor dela.
(c)~$2^{-n} < 0.01$ quando $n > \log_2 100 = 6.64$: após cerca de sete meias-vidas.
\end{solution}

\begin{solution}{exo:b1:radiometric-methods:2}
(a)~Radiocarbono no carvão. (b)~Urânio--chumbo no zircão, ou argônio--argônio
(potássio--argônio) na sanidina, um \omterm{def:b1:rock-forming-minerals:mineral}{mineral} rico em potássio. (c)~Urânio--chumbo no zircão,
que sobrevive a uma longa história; rubídio--estrôncio em rochas totais é possível;
potássio--argônio correria o risco de idades rebaixadas pelo argônio perdido em aquecimentos
posteriores. (d)~Radiocarbono no carbonato da concha, perto do limite do método (cerca de
\num{70000} anos). (e)~Nenhum diretamente: sua idade é contida por rochas vulcânicas datadas
acima e abaixo, ou obtida correlacionando seus fósseis com uma seção datada
(\cref{ch:b1:stratigraphic-scale}).
\end{solution}

\begin{solution}{exo:b1:radiometric-methods:3}
(a)~$t = -8033 \ln 0.25 = 8033 \times 1.3863 = \num{11136}$ anos BP.
(b)~A vida média é $5700/\ln 2 = \num{8223}$ anos: $t = 8223 \times 1.3863 =
\num{11400}$ anos. (c)~Ambos supõem que a proporção de carbono-14 no carbono vivo era a mesma
então que no padrão; não era. A idade civil é lida em uma \omterm{def:b1:radiometric-methods:radiocarbon}{curva de calibração} como IntCal20,
como uma faixa em anos cal BP.
\end{solution}

\begin{solution}{exo:b1:radiometric-methods:4}
(a)~Cada relógio começa quando seu \omterm{def:b1:rock-forming-minerals:mineral}{mineral} esfria abaixo de sua \omterm{def:b1:radiometric-methods:closure}{temperatura de fechamento}:
zircão (acima de \qty{750}{\celsius}) fecha primeiro, a hornblenda (cerca de
\qty{530}{\celsius}) em seguida, a biotita (cerca de \qty{280}{\celsius}) por último, de modo
que as idades diminuem nessa ordem à medida que o granito esfriou. (b)~$(530 - 280)/(410 - 398) =
250/12 = \qty{21}{\celsius}$ por milhão de anos. (c)~O zircão fechou acima de
\qty{750}{\celsius}, então o granito esfriou ao menos $750 - 530 =
\qty{220}{\celsius}$ em \qty{10}{\Myr}: ao menos \qty{22}{\celsius} por milhão de anos, um
limite inferior.
\end{solution}

\begin{solution}{exo:b1:radiometric-methods:5}
(a)~$\bar x = 7.5$, $\bar y = 0.809525$, $S_{xx} = 81.0$, $S_{xy} = 1.13965$;
$m = 1.13965/81.0 = 0.0140698$, $b = 0.809525 - 7.5 \times 0.0140698 = 0.70400$.
(b)~$t = \ln(1.0140698)/\num{1.3972e-11} = \qty{1000.0}{\Ma}$; \omterm{def:b1:radiometric-methods:isochron}{razão inicial}
$0.7040$. (c)~Os resíduos são $+1.6$, $-2.5$, $+0.7$ e $+0.2$ (em unidades de
\num{e-4}): $s = \sqrt{\num{9.34e-8}/2} = 0.00022$, $\sigma_m = 0.00022/\sqrt{81}
= \num{2.40e-5}$ e $\sigma_t = \num{2.40e-5}/(\num{1.3972e-11} \times 1.0141) =
\qty{1.7}{\Myr}$: \qty[uncertainty-mode=separate]{1000.0(17)}{\Ma}. (d)~$0.7040$ está próximo
do valor do \omterm{def:b1:layered-earth:layers}{manto} $0.7025$ e bem abaixo do da \omterm{def:b1:layered-earth:layers}{crosta} $0.710$: uma fonte sobretudo no \omterm{def:b1:layered-earth:layers}{manto}.
\end{solution}

\begin{solution}{exo:b1:radiometric-methods:6}
(a)~$0.100/39.098 \times \num{1.17e-4} = \qty{2.992e-7}{mol}$ de potássio-40 por grama.
(b)~${}^{40}\mathrm{Ar}^*/{}^{40}\mathrm{K} = \num{1.169e-9}/\num{2.992e-7}
= 0.0039065$; $\lambda/\lambda_e = 5.543/0.581 = 9.5404$; $t = \ln(1 + 9.5404
\times 0.0039065)/\num{5.543e-10} = \ln(1.037269)/\num{5.543e-10} =
\qty{66.01}{\Ma}$. (c)~O argônio radiogênico seria $0.995 \times 1.169$ nanomols, dando
\qty{65.69}{\Ma}: a idade medida estaria \qty{0.32}{\Myr} velha demais, dez vezes os
\num{32000} anos dentro dos quais o impacto e a extinção no fim do Cretáceo são agora mostrados
coincidir. O excesso de argônio preocupa sobretudo as rochas jovens e as rochas totais; um
\omterm{def:b1:rock-forming-minerals:mineral}{mineral} como a sanidina não deveria incorporar argônio ao cristalizar, e um espectro de idades
argônio--argônio o revelaria.
\end{solution}

\begin{solution}{exo:b1:radiometric-methods:7}
(a)~Cerca de \num{2000} cal BP: a curva passa por \num{2070} anos BP ali. (b)~De cerca de
\num{2495} a \num{2699} cal BP, duzentos anos: a curva é quase plana entre \num{2500} e
\num{2700} cal BP, de modo que toda idade civil desse intervalo tem quase a mesma \omterm{def:b1:radiometric-methods:radiocarbon}{idade
radiocarbono}, e uma medida mais precisa não estreitaria a faixa. (c)~A \num{2000} cal BP a \omterm{def:b1:radiometric-methods:radiocarbon}{idade
radiocarbono} é \num{70} anos \emph{mais velha} que a idade civil; a \num{8000} cal BP,
\num{822} anos mais jovem (\qty{10}{\percent}); a \num{12000} cal BP \num{1694} anos mais jovem
(\qty{14}{\percent}). A diferença não é proporcional à idade: ela chega até a mudar de sinal.
Ela não pode vir só da meia-vida, que daria um \qty{2.4}{\percent} constante.
\end{solution}

\begin{solution}{exo:b1:radiometric-methods:8}
(a)~A camada de cinza abaixo (\qty{94.50}{\Ma}) dá uma idade máxima, o derrame de lava acima
(\qty{93.80}{\Ma}) uma mínima: o limite fica contido entre elas. O dique também é uma mínima,
mas mais fraca, mais jovem que a lava. O zircão detrítico é uma máxima apenas para o arenito, e
mais fraca que a camada de cinza, que fica abaixo do arenito. (b)~$f = (12 - 4)/(20 - 4) = 0.5$;
$t = 94.50 - 0.5 \times 0.70 =
\qty{94.15}{\Ma}$; $\sigma_t = \sqrt{0.25 \times 0.10^2 + 0.25 \times 0.12^2} =
\qty{0.08}{\Myr}$: \qty[uncertainty-mode=separate]{94.15(8)}{\Ma}. (c)~Seus
grãos vieram de uma rocha de ao menos \qty{95.3}{\Ma}, erodida em algum lugar da área-fonte; isso
mostra que o arenito é mais jovem que \qty{95.3}{\Ma}, o que a camada de cinza abaixo já prova
com mais rigor.
\end{solution}

\begin{solution}{exo:b1:radiometric-methods:9}
(a)~Grão~A: $\ln(1.08401)/\num{1.55125e-10} = \qty{520.0}{\Ma}$ e
$\ln(1.6688)/\num{9.8485e-10} = \qty{520.0}{\Ma}$. Grão~B: \qty{368.4}{\Ma} e
\qty{389.9}{\Ma}. (b)~O grão~A é concordante. (c)~$0.05881/0.08401 = 0.7000$ e
$0.4682/0.6688 = 0.7001$. A idade é $\ln(1 + R)/\lambda$, não proporcional a $R$: para a
pequena razão do urânio-238 $\ln(1+R) \approx R$ e a idade cai quase \qty{30}{\percent}
(para \qty{368}{\Ma}); para a razão maior do urânio-235 o logaritmo cai menos, e a idade só
cai para \qty{390}{\Ma}. Ambas são jovens demais, e elas discordam: o grão é discordante.
\end{solution}

\begin{solution}{exo:b1:radiometric-methods:10}
(a)~Com $u = \lambda t$, $D^*/P = e^u - 1 = u + u^2/2 + \dots$, de modo que $(D^*/P)/\lambda
= t\,(1 + u/2 + \dots)$: a aproximação superestima a idade por um erro relativo próximo de
$u/2$. (b)~Rubídio-87 a \qty{310}{\Ma}: $u = 0.00433$, erro \qty{0.22}{\percent}
(\qty{0.7}{\Myr}); potássio-40 a \qty{296}{\Ma}: $u =
0.164$, erro \qty{8.7}{\percent} (\qty{26}{\Myr}; $u/2$ dá \qty{8.2}{\percent});
urânio-238 a \qty{4.4}{\Ga}: $u = 0.683$, erro \qty{43}{\percent} --- a aproximação e até seu
erro estimado ($u/2 = \qty{34}{\percent}$) falham. (c)~$R = 0.004341$, $\sigma_R = \num{4.34e-6}$; pela \cref{prop:b1:radiometric-methods:precision}, $\sigma_t =
\num{4.34e-6}/(\num{1.3972e-11} \times 1.0043) = \qty{0.31}{\Myr}$.
\end{solution}

\begin{solution}{exo:b1:radiometric-methods:11}
(a)~Seja ${}^{40}\mathrm{K}$ o potássio-40 de hoje. Em $t_2$ antes do presente havia
${}^{40}\mathrm{K}\,e^{\lambda t_2}$ átomos, e o argônio produzido entre $t_1$ e $t_2$ foi,
pela \cref{prop:b1:radiometric-methods:kar} aplicada a esse intervalo, $(\lambda_e/\lambda)\,{}^{40}\mathrm{K}\,e^{\lambda
t_2}\bigl(e^{\lambda (t_1 - t_2)} - 1\bigr) = (\lambda_e/\lambda)\,
{}^{40}\mathrm{K}\,\bigl(e^{\lambda t_1} - e^{\lambda t_2}\bigr)$. Uma fração $1 - f$ dele
permaneceu. Desde $t_2$, o \omterm{def:b1:rock-forming-minerals:mineral}{mineral} fechado produziu $(\lambda_e/\lambda)\,{}^{40}\mathrm{K}\,(e^{\lambda t_2} - 1)$
a mais. Somando e multiplicando por $\lambda/(\lambda_e\,{}^{40}\mathrm{K})$ obtém-se a expressão.
(b)~A idade aparente é $\ln(1 + \text{expressão})/\lambda$: \qty{400.0}{\Myr} para $f = 0$,
\qty{342.3}{\Myr} para $0.3$, \qty{282.7}{\Myr} para $0.6$ e \qty{200.0}{\Myr} para $1$.
(c)~Toda idade entre \qty{200}{\Myr} e \qty{400}{\Myr} pode ser produzida por alguma perda
parcial, e um único número não pode dizer se ele registra um evento real naquele tempo ou uma
mistura; um espectro de idades, ou outros relógios na mesma rocha, são necessários.
\end{solution}

\begin{solution}{exo:b1:radiometric-methods:12}
(a)~A idade é proporcional à vida média usada: $10306 \times 8223/8033 =
\num{10550}$ anos. (b)~$10550 - 10306 = 244$ anos, cerca de um sétimo dos
\num{1694}. (c)~Se a atividade inicial era $k A_0$, a idade verdadeira é $t = \tau
\ln(k A_0/A) = \num{10550} + \tau \ln k$; com $t = \num{12000}$, $\ln k =
1450/8223 = 0.176$ e $k = 1.19$: cerca de \qty{19}{\percent} mais carbono-14 que
no padrão.
\end{solution}

\begin{solution}{pb:b1:radiometric-methods:1}
\begin{enumerate}
  \item Os pontos sobem de modo constante com $x$; seus desvios de uma reta são de no máximo
        duas unidades na quarta casa decimal, contra uma amplitude de $0.04$ em $y$: estão sobre
        uma reta dentro de uma pequena dispersão.
  \item $\bar x = 4.90333$, $\bar y = 0.729288$, $S_{xx} = 62.1253$, $S_{xy} =
        0.269335$; $m = 0.269335/62.1253 = 0.0043354$.
  \item $b = 0.729288 - 0.0043354 \times 4.90333 = 0.70803$: a $^{87}\mathrm{Sr}/^{86}\mathrm{Sr}$
        do granito quando cristalizou, compartilhada por todas as suas partes.
  \item $t = \ln(1.0043354)/\num{1.3972e-11} = \qty{309.6}{\Ma}$.
  \item Os resíduos são $+9.1$, $-17.2$, $+6.9$, $-6.4$, $+15.1$ e $-7.5$ (em unidades de
        \num{e-5}); $s = 0.000137$, $\sigma_m = 0.000137/\sqrt{62.1253} =
        \num{1.74e-5}$, $\sigma_t = \num{1.74e-5}/(\num{1.3972e-11} \times 1.0043)
        = \qty{1.2}{\Myr}$: \qty[uncertainty-mode=separate]{309.6(12)}{\Ma}, e
        $\sigma_b = 0.00010$.
  \item $0.7080$ está mais perto do $0.710$ da \omterm{def:b1:layered-earth:layers}{crosta} do que do $0.7025$ do \omterm{def:b1:layered-earth:layers}{manto}. Para uma
        mistura, $(0.70803 - 0.7025)/(0.710 - 0.7025) = 0.74$: cerca de três quartos do
        estrôncio vieram de \omterm{def:b1:layered-earth:layers}{crosta} fundida.
  \item $0.0740/39.098 \times \num{1.17e-4} = \qty{2.214e-7}{mol}$ por grama, vezes
        $\num{6.022e23}$: \num{1.33e17} átomos.
  \item $^{40}\mathrm{Ar}^*/^{40}\mathrm{K} = \num{4.139e-9}/\num{2.2144e-7} =
        0.018691$; $1 + 9.5404 \times 0.018691 = 1.17832$; $t =
        \ln(1.17832)/\num{5.543e-10} = \qty{296.0}{\Ma}$.
  \item $^{40}\mathrm{K} = \qty{3.292e-8}{mol.g^{-1}}$, razão $0.019239$,
        $\ln(1.18355)/\num{5.543e-10} = \qty{304.0}{\Ma}$.
  \item $\ln(1.018691)/\num{5.543e-10} = \qty{33.4}{\Ma}$, nove vezes jovem demais, porque só
        cerca de um decaimento em dez dá argônio. O cálcio-40 é o isótopo comum do cálcio,
        abundante nas rochas: o pouco produzido por decaimento não pode ser distinguido do cálcio
        que os \omterm{def:b1:rock-forming-minerals:mineral}{minerais} continham desde o início.
  \item O excesso de argônio desloca idades em quantidades da ordem de um milhão de anos, o que
        importa para rochas jovens; aqui seria bem abaixo de \qty{1}{\percent} da idade. E
        cristais de \omterm{def:b1:rock-forming-minerals:groups}{mica} e de \omterm{def:b1:rock-forming-minerals:groups}{anfibólio} quase não incorporam argônio ao se formar.
  \item Ou o granito esfriou devagar e as \omterm{def:b1:rock-forming-minerals:groups}{micas} e o \omterm{def:b1:rock-forming-minerals:groups}{anfibólio} fecharam-se depois da cristalização
        (\omterm{def:b1:radiometric-methods:closure}{idades de resfriamento}), ou um aquecimento posterior expulsou parte de seu argônio. Um
        patamar plano em um espectro de idades argônio--argônio apoia uma \omterm{def:b1:radiometric-methods:closure}{idade de resfriamento};
        um espectro sem patamar alerta que os \omterm{def:b1:rock-forming-minerals:mineral}{minerais} perderam argônio. A ordem das idades, a
        hornblenda antes da biotita como preveem suas \omterm{def:b1:radiometric-methods:closure}{temperaturas de fechamento}, também aponta
        para o resfriamento.
  \item $\ln(1.049361)/\num{1.55125e-10} = \qty{310.6}{\Ma}$ e
        $\ln(1.35784)/\num{9.8485e-10} = \qty{310.6}{\Ma}$: concordantes.
  \item Zircão \qty{310.6}{\Ma}, fechamento acima de \qty{750}{\celsius}: cristalização; a
        \omterm{def:b1:radiometric-methods:isochron}{isócrona} de rocha total, \qty[uncertainty-mode=separate]{309.6(12)}{\Ma}, também a data,
        e concorda dentro de sua incerteza; hornblenda \qty{304.0}{\Ma}: resfriamento através de
        cerca de \qty{530}{\celsius}; biotita \qty{296.0}{\Ma}: resfriamento através de cerca de
        \qty{280}{\celsius}. O gráfico desce de acima de \qty{750}{\celsius} em
        \qty{310.6}{\Ma} até \qty{280}{\celsius} em \qty{296.0}{\Ma}.
  \item $(530 - 280)/(304.02 - 296.03) = 250/7.99 = \qty{31}{\celsius}$ por milhão de anos.
  \item $(750 - 530)/(310.60 - 304.02) = 220/6.58 = \qty{33}{\celsius}$ por milhão de anos ao
        menos, já que o zircão fechou acima de \qty{750}{\celsius}: o resfriamento foi ao menos
        tão rápido no início.
  \item A diferença de temperaturas, $250 \pm \sqrt{40^2 + 40^2} = 250 \pm
        57$, dá
        \qtyrange{24}{38}{\celsius} por milhão de anos.
  \item $(550 - 350)/7.99 = \qty{25}{\celsius}$ por milhão de anos, dentro da faixa da questão~17:
        a incerteza vem das \omterm{def:b1:radiometric-methods:closure}{temperaturas de fechamento}, não das idades.
  \item $(280 - 10)/30 = \qty{9.0}{km}$.
  \item $\ln(1.039108)/\num{1.55125e-10} = \qty{247.30}{\Ma}$, e o mesmo pela razão do
        urânio-235: concordante.
  \item Os zircões detríticos vêm do granito, erodido na \omterm{def:b1:cross-sections:unconformity}{discordância}: sua idade é a idade de
        cristalização do granito. Eles mostram apenas que o arenito é mais jovem que
        \qty{310}{\Ma}: uma idade máxima, muito mais fraca que a camada de cinza.
  \item $(58 - 18)/(247.30 - 245.30) = \qty{20.0}{m}$ por milhão de anos; a base é
        $18/20.0 = \qty{0.90}{\Myr}$ mais velha que a cinza inferior: \qty{248.20}{\Ma}.
  \item De \qty{309.6}{\Ma} (\omterm{def:b1:radiometric-methods:isochron}{isócrona}) ou \qty{310.6}{\Ma} (zircão) até \qty{248.2}{\Ma}: cerca
        de \qty{62}{\Myr}.
  \item \qty{9.0}{km} em $296.0 - 248.2 = \qty{47.8}{\Myr}$: \qty{0.19}{km} por milhão de anos,
        \qty{0.19}{mm} por ano. Isso supõe que a biotita datada veio logo abaixo da \omterm{def:b1:cross-sections:unconformity}{discordância},
        de modo que toda a rocha acima dela no granito e no xisto foi removida quando a areia
        chegou, e que o geoterma permaneceu constante.
  \item O granito esfriou a cerca de \qty{31}{\celsius} por milhão de anos entre os fechamentos
        da hornblenda e da biotita.
\end{enumerate}
\end{solution}
